\documentclass[12pt]{article}
\usepackage[utf8]{inputenc}
\usepackage[T1]{fontenc}
\usepackage{hyperref}
\title{Amba Language Resources for K-12 Curriculum Development}
\author{Compiled by Claude Sonnet 5 (claude-sonnet-5) for Andr\'as Kornai}
\date{}
\begin{document}
\maketitle
\section{Introduction}
Amba (ISO 639-3 \texttt{utp}; also Aba, Nembao, Nebao) is an Oceanic language of Utupua in the Santa Cruz group, Solomon Islands (Glottolog amba1266). Wikipedia gives 590 native speakers, from a 1999 count. The Endangered Languages site rates it ``severely endangered'', but with low certainty and on older data. A 2022 expedition report says Utupua's languages are highly endangered and almost completely undocumented. The spelling of the name has not been standardized.
\section{Primary Educational Resources}
I found no primers, textbooks or school programmes for Amba.
\section{Linguistic and Reference Materials}
\begin{itemize}
\item Amba language (Solomon Islands), Wikipedia: \url{https://en.wikipedia.org/wiki/Amba_language_(Solomon_Islands)}
\item Endangered Languages Project entry: \url{https://www.endangeredlanguages.com/lang/3871?hl=es}
\item Verbix page on Amba: \url{https://docs.verbix.com/Languages/Amba}
\end{itemize}
\section{Religious and Community Resources}
None found.
\section{Digital and Audio Resources}
None found.
\section{Curriculum Development Context}
The language is essentially undocumented, so no curriculum can currently draw on existing native-language materials. Documentation and orthography work with the Utupua community would have to come first.
\end{document}